\chapter{Ориентированная спектральная самоэнергия притока}
\label{ch:v10-inflow-spectral-self-energy-running-parent-origin}

\section{Проверяемая гипотеза}

Предыдущий гейт установил, что умножение одинаковых ячеек создаёт только
экстенсивный ход: после деления на объём интенсивная $\beta$-функция
исчезает. Поэтому настоящий гейт проверяет более сильный механизм. При
геометрическом росте должен меняться не только размер резервуара, но и
локальный спектральный отклик одной входящей степени свободы.

Проверка разделена на две части. Сначала строится минимальная точная
архитектура такого отклика. Затем отдельно отмечается, выведено ли её
вложение в уже существующий $43$-мерный носитель программы.

\section{Взаимно обратный спектральный резервуар}

Пусть $\zeta=\log a$ --- геометрическая координата роста. Рассмотрим
двухмодовый положительный оператор
\begin{equation}
 K_{\rm res}(\zeta)
 =\begin{pmatrix}e^{-\zeta}&0\\0&e^{\zeta}\end{pmatrix},
 \qquad
 \det K_{\rm res}=1,
 \label{eq:v10-reciprocal-reservoir}
\end{equation}
и его точный обратный оператор
\begin{equation}
 G_{\rm res}(\zeta)=K_{\rm res}^{-1}(\zeta)
 =\begin{pmatrix}e^{\zeta}&0\\0&e^{-\zeta}\end{pmatrix}.
 \label{eq:v10-reciprocal-propagator}
\end{equation}
Определитель остаётся единичным: сама взаимно обратная пара не создаёт
аномалию и не меняет суммарное спектральное число.

Две ориентации выделяются дополнительными проекторами
\begin{equation}
 P_{\rm in}=\begin{pmatrix}1&0\\0&0\end{pmatrix},
 \qquad
 P_{\rm out}=\begin{pmatrix}0&0\\0&1\end{pmatrix},
 \qquad P_{\rm in}+P_{\rm out}=I_2.
 \label{eq:v10-orientation-projectors}
\end{equation}
Они имеют ранг один. Геометрическая стрелка рождения ячеек условно выбирает
$P_{\rm in}$; противоположный проектор служит точным контрольным каналом.

\section{Самоэнергия как ориентированный спектральный отклик}

Для единичных векторов $b_{\rm in}=(1,0)^T$ и
$b_{\rm out}=(0,1)^T$ соответствующие дополнения Шура равны
\begin{equation}
 \Sigma_{\rm in}(\zeta)
 =b_{\rm in}^{T}G_{\rm res}b_{\rm in}=e^{\zeta},
 \qquad
 \Sigma_{\rm out}(\zeta)
 =b_{\rm out}^{T}G_{\rm res}b_{\rm out}=e^{-\zeta}.
 \label{eq:v10-oriented-self-energies}
\end{equation}
Следовательно,
\begin{equation}
 \Sigma_{\rm in}\Sigma_{\rm out}=1,
 \qquad
 \frac{d\Sigma_{\rm in}}{d\zeta}=\Sigma_{\rm in},
 \qquad
 \frac{d\Sigma_{\rm out}}{d\zeta}=-\Sigma_{\rm out}.
 \label{eq:v10-oriented-self-energy-beta}
\end{equation}
Это уже интенсивный ход: он относится к локальному спектральному отклику,
а не к числу одинаковых копий.

Ориентация существенна. Симметричный вектор
$b_{\rm sym}=2^{-1/2}(1,1)^T$ даёт
\begin{equation}
 \Sigma_{\rm sym}(\zeta)=\cosh\zeta,
 \qquad
 \left.\frac{d\Sigma_{\rm sym}}{d\zeta}\right|_{\zeta=0}=0.
 \label{eq:v10-symmetric-self-energy-zero-beta}
\end{equation}
Фиксированный локальный спектр также имеет нулевой ход. Поэтому ненулевая
производная не следует из одного двухмодового удвоения: её создаёт совместное
наличие взаимно обратного спектра и направленной связи.

\section{Ненулевой свидетель следового отклика}

Подставим входящую самоэнергию в нормированное действие предыдущего гейта:
\begin{equation}
 \Gamma_{\rm in}(q,\zeta)
 =\log(1+q^2+e^{\zeta})
  -\log(1+q^2)-\log(1+e^{\zeta}).
 \label{eq:v10-running-inflow-action}
\end{equation}
Его геометрическая плотность отклика
\begin{equation}
 \mathcal A(q,\zeta)
 :=\frac{\partial\Gamma_{\rm in}}{\partial\zeta}
 =\frac{e^{\zeta}}{1+q^2+e^{\zeta}}
  -\frac{e^{\zeta}}{1+e^{\zeta}}
 \label{eq:v10-running-trace-response}
\end{equation}
не равна нулю. В рациональной контрольной точке
\begin{equation}
 \mathcal A(1,0)=-\frac16.
 \label{eq:v10-running-trace-response-witness}
\end{equation}
Тем самым построен точный ненулевой интенсивный свидетель, которого не было
при чистом объёмном умножении ячеек.

\begin{theorem}[Ориентированная самоэнергия имеет ненулевой ход]
Взаимно обратный резервуар с $\det K_{\rm res}=1$ сам по себе не создаёт
масштаб. После выбора входящего ранга один его дополнение Шура равно
$e^{\zeta}$, имеет ненулевую интенсивную $\beta$-функцию и создаёт
$\mathcal A(1,0)=-1/6$. Противоположная ориентация имеет обратный ход, а
симметричная связь имеет нулевую первую производную в $\zeta=0$.
\end{theorem}

\begin{nogocriterion}[Абстрактный резервуар ещё не является физическим]
Двухмодовый оператор и проектор ориентации образуют согласованную
конечномерную архитектуру, но не доказывают, что именно этот блок содержится
в $43$-мерном концевом носителе, согласован с шестимерным KMS-сектором и
порождается тем же родительским функционалом. Без такого типизированного
вложения ненулевой отклик остаётся условным представителем, а не выводом
абсолютного масштаба.
\end{nogocriterion}

\begin{protocol}[Статус спектрального гейта]\begin{enumerate}
\item взаимно обратный положительный резервуар: \textbf{построен};
\item входящая и выходящая ориентации: \textbf{ранг $1+1$};
\item интенсивный спектральный ход: \textbf{ненулевой};
\item точный следовой свидетель: \textbf{$-1/6$};
\item архитектура ориентированного отклика: \textbf{$7/7$};
\item происхождение геометрии и стрелки роста: \textbf{$2/2$ условно
унаследовано};
\item типизированное физическое вложение резервуара: \textbf{$0/1$};
\item полный физический родитель абсолютного масштаба: \textbf{$0/1$};
\item следующий гейт: \textbf{типизированное вложение спектральной
самоэнергии в $43$-мерный носитель}.
\end{enumerate}\end{protocol}

Машинный сертификат сохранён в
\path{s2t_v10_inflow_spectral_self_energy_running_parent_origin_gate_results.json}.

\begin{thebibliography}{9}
\bibitem{v10-running-keldysh}
L.~M.~Sieberer, M.~Buchhold, S.~Diehl,
\emph{Keldysh Field Theory for Driven Open Quantum Systems},
Reports on Progress in Physics 79 (2016), 096001.

\bibitem{v10-running-feshbach}
H.~Feshbach,
\emph{A Unified Theory of Nuclear Reactions. II},
Annals of Physics 19 (1962), 287--313.
\end{thebibliography}